OCAML is a compiled, functional programming language and a descendant of the Meta Language programming language.
In contrast to other functional programming languages (such as Haskell), it is actually used in the wild by more than two people.

\inputcodewithfilename{ocaml}{}{code/ocaml/00_basics.ml}

OCAML is dynamically typed (i.e. types don't \textit{have} to be annotated), but they can be, if you so choose.
As with Python and all other languages where type annotation is optional, you are strongly encouraged to annotate types, especially for more complex expressions,
as that can massively reduce the likeliness of errors.

The first graded assignment's first 300 ish lines also give you a very neat overview over how OCAML works,
so absolutely read through those.

For \texttt{Printf}, the format string is equivalent to that of \texttt{C},
a solid cheat sheet can be found \hlhref{https://alvinalexander.com/programming/printf-format-cheat-sheet/}{here}
